% @BAKE pdflatex $@
\documentclass[12pt]{article}

% Encoding and language
\usepackage[T1]{fontenc}
\usepackage[utf8]{inputenc} % Remove if using LuaLaTeX or XeLaTeX
\usepackage[english]{babel}

% Better typography
\usepackage{lmodern}
\usepackage{microtype}

% Page layout
\usepackage[a4paper,margin=10mm]{geometry}

\usepackage{pgfpages}
\pgfpagesuselayout{4 on 1}[a4paper,border shrink=5mm]

\usepackage{anyfontsize}

\title{Dino Side-Quest}
\date{\today}

\begin{document}

% Title page
\begin{titlepage}
    \fontsize{22pt}{22pt}\selectfont
    \maketitle
    Én \${NÉV} vállalom,
    hogy ezt a documentumot az alá írást követő 7 napban elolvasom.

    \vspace{1cm}
    \dotfill

    \vspace{1cm}
    \dotfill

    \vspace{1cm}
    \dotfill

    \vfill
\end{titlepage}

\section*{The Eiffel Tower}
\fontsize{25pt}{22pt}\selectfont

\bigskip

When I was working on the construction of the Eiffel Tower, yes,
those were the days. And yet I never knew how happy I was.

The construction of the Eiffel Tower was a most beautiful and
important enterprise. Today you can’t imagine what it meant. What
is presently the Eiffel Tower has very little to do with the reality of that
time. Since then the dimensions have changed; it seems as if the tower
has contracted. I walk beneath it, raise my eyes, look. But
fail to rec-
ognize the world where I spent the most beautiful days of my life.

Tourists enter the elevator, go up to the first gallery, the second, ex-
claim, laugh, take photographs, shoot color films. Yet they don’t
know, the poor devils, and they will never know.

In the guidebooks you read that the Eiffel Tower stands 300 me-
ters tall; the radio antenna adds another 20 meters to this height. The
newspapers of the period, even before the workers began, gave the
same figure. And to the public 300 meters seemed utter madness.

Yet 300 is really nothing. I used to work at the Runtiron offices,
near Neuilly. I was a skilled mechanical worker. One evening, asI was
setting out for home, a gentleman about forty wearing a top hat
stopped me on the street.

“Are you Monsieur André Lejeune?” he asked me.

“Precisely,” I answered. “And who are you, Monsieur?”

“Tam the engineer Gustave Eiffel, and I would like to make you
an offer. But first Imust show you something. Here is my carriage.”

I got into the engineer’s carriage, and he took me to a huge shed
constructed in a field on the outskirts of the city. Here there were
about thirty young men who worked in silence at large drawing ta-
bles. None of them thought me worthy of a glance.

The engineer led me to the back of the building where a drawing
of atower about two meters long was tacked on the wall. “I shall con-
struct this tower you see here for Paris, for France, for the world. It
will be made of iron, and it will be the tallest tower in the world.”

“How high will it be?” I asked.

“The official project anticipates a height of 300 meters. This is the
figure we stipulated in our contract with the government, so they
wouldn’t be frightened. The tower, however, will be much taller.”

“Four hundred meters?”

“Young man, believe me, I cannot discuss it now. You must have
patience. Nonetheless, it will be a marvelous undertaking; to partici-
pate in it will be an honor. And I have come looking for you because I
have been informed that you are an expert worker. How much do you
make at Runtiron?” I told him my salary. “If you come and work for
me,” said the engineer, brusquely shifting to a more familiar tone,
“TI pay you three times as much.” I accepted the offer.

Then the engineer said in a low voice, “I forgot one thing, dear
André. I attach great importance to yourjoining us, but first you must
make a promise.”

“T hope it will be nothing less than honorable,” I hazarded, some-
what taken aback by that air of mystery.

“Itis the secret,” he said.

“What secret?”

“Will you give me your word of honor that you will not speak to
anyone, not even to the members of your family, about anything that
concerns our work? Do you promise that you will not tell a living soul
what you will do or how you will do it? That you will not reveal num-
bers, measurements, data, or ciphers? Think about it very carefully
before you shake my hand: one day this secret will weigh heavy on
you.”

On the table lay a printed form, the contract for the job, and this
document contained the pledge to maintain the secret. I signed.

There were hundreds, perhaps thousands, of workers on the site.

Not only did I not recognize any of them, but after we began I never
saw them, since we worked in groups without a break in the conti-
nuity and there were three shifts in every twenty-four-hour period.

After the cement foundation had been laid, we mechanics began
to mount the steel girders. From the beginning, there was very little
talk among us, perhaps because we had sworn to keep the secret. But
from some comments I had caught here and there, I got the idea that
my fellow workers had accepted the job only because of the excep-
tional salary. No one, or nearly no one, believed that the tower would
ever be finished. They considered it a folly, beyond the capabilities of
human strength.

Nevertheless, with the four gigantic feet planted in the earth, the
iron framework rose into view. On the other side of the fence, at the
edges of the vast work site, a crowd stationed itself day and night to
gaze at us, as we maneuvered up above like ants, hanging from the
web.

The arches of the pedestal were successfully welded, the four ver-
tical columns rose nearly perpendicular and then were joined to form
a single column which gradually grew thinner. In eight months the
tower had reached a height of 100 meters, and the workers were given
an outdoor banquet, at a restaurant on the Seine.

Ino longer heard any words of mistrust. In fact, a strange enthu-
siasm pervaded workers, foremen, technicians, and engineers, as if
we were all on the eve of an extraordinary event. One morning, early
in October, we found ourselves immersed ina fog.

It was thought that a blanket of low clouds had come to a stand-
still over Paris, but this wasn’t so. All around the sky was clear. “Look
at that tube,” said Claude Gallumet, the smallest and quickest worker
in my group, who had become my friend. Some sort of whitish
smoke issued from a huge rubber tube fixed to the iron framework.

There were four tubes, one for each corner of the tower. They pro-
duced a dense vapor, gradually forming a cloud that remained en-
tirely motionless; amid this great cotton-wool umbrella we contin-
ued to work. But why was it there? Because of the secret?
When we reached 200 meters, the builders gave us another ban-
quet. Even the newspapers mentioned the great occasion. Yet the
crowd was no longer stationed around the work site; that ridiculous
cap of cloud hid us completely from their gazes. And the newspapers
praised the device: that condensation of vapor, they said, kept the
workers who were hoisted up on the aerial structures from seeing the
abyss beneath them, and thus it prevented vertigo. What patent non-
sense! First of all, every one of us was quite accustomed to the height
at this point; in cases of vertigo, moreover, no mishaps would have
occurred because each of us wore a solid leather belt which was se-
curely fastened to the surrounding structures by a cable.

The tower continued to rise; we reached 250, 280, 300 meters.

Almost two years had now passed. Were we at the end of our adven-
ture? One evening we were assembled under the great cross of the
base, and the engineer Eiffel spoke to us. Our undertaking, he said,
had come to an end, we had given proof of our perseverance, skill,
and courage, and in fact, the building firm was awarding us a special
bonus. Whoever wanted to return to his home could do so. Yet Eiffel
hoped that there might be some volunteers who were willing to con-
tinue with him. Continue what? The engineer could not explain what
he meant, the workers must put their trust in him, it would be worth
the pain.

I stayed on with many others. It was a kind of mad conspiracy
which no stranger suspected because each one of us remained faithful
to the secret more than ever.

So, at a height of 300 meters, rather than top off the structure
with a spire, we mounted more steel girders one on top of the other
toward the zenith. One iron beam was attached to another, one girder
to another, with countless bolts and the din of hammers, the entire
cloud vibrated with them like a sound box. We were in flight.

Until, by dint of the ascent, we had left behind the group of
clouds, which remained beneath us, and although the people of Paris
were still unable to see us because of that vaporous screen, in reality
we expatiated in the pure and limpid air of the upper reaches. On cer-
tain windy mornings we could see the distant Alps covered with
snow.

By this time the tower had reached such a height that traveling up
and down it was taking us more than half the working hours. The el-
evators weren’t useful any more. From day to day the time available
for productive work decreased. The day would come when we would
have to start our descent immediately after we arrived at the top.

And then the tower would not rise any more, not even a meter higher.

So we decided to install some of our sheds up there, between the
iron beams, like nests, which could not be seen from the city because
they were hidden by the mist of artificial clouds. There we slept, ate,
and at night played cards, when the great chorus of illusions and vic-
tories was not in harmony. We took turns going down to the city, but
only on holidays.

It was in that period that you slowly perceived the marvelous
truth, the motive, that is, for the secret. And we no longer felt like
mechanical workers. We were pioneers, explorers, we were heroes,
saints. You slowly began to realize that the construction of the Eiffel
Tower would never end, now you understood why the engineer had
wanted the sesquipedalian pedestal, those four cyclopean iron claws
which seemed absolutely exaggerated. The construction would
never be finished, and for eternity the Eiffel Tower would continue to
rise toward the sky, advancing beyond the clouds, the tempests, the
peaks of Gaurisangar. As long as God gave us strength, we would
continue to bolt one steel beam on top of the other, always higher and
higher, and after us our sons would continue the work, and no one in
the flat city of Paris would know, the dreary world would never
know.

Of course, sooner or later, they would lose their patience down
there, protests and petitions would be made to Parliament, people
would say, “Why on earth haven't they stopped building that blasted
tower? They’ve already reached the anticipated 300 meters, so they
should have finished off the thing.” Yet we would find pretexts, in
fact, we would manage to put some of our men in Parliament or
among the ministers, we would hush up the affair, the people of the
lower world would resign themselves, and we would be in supreme
exile, always farther and farther up in the sky.

The sound of gunfire was heard down below, beyond the white
cloud. We went down a considerable distance, passed through the
fog, and looked out from its lower boundary with a telescope. To-
ward our work site, in a concentric formation and with a threatening
attitude, advanced the gendarmes from the royal guard, the police,
the bureaus of inspection, the garrisons, the battalions, the army, and
the navy. May the Devil flay and devour them all!
They sent up a courier with an ultimatum: if you bastards don’t
surrender and come down immediately, in six hours we'll open fire
with rifles, machine guns, and light artillery, and you'll be done for.

So a filthy Judas had betrayed us. The son of the engineer Eiffel,
whose great father was dead and buried for some time now, turned
pale as a ghost. How could we fight? Thinking of our families, we
surrendered.

They took apart the poem we had raised to the heavens, ampu-
tated the spire at a height of 300 meters, and thrust on it that ugly little
cap which you still see today.

The cloud that hid us no longer exists; because of this cloud, in
fact, there will be a hearing in a Parisian court. The aborted tower was
painted entirely gray, and from it hang long banners which wave in
the sun. Today is the inaugural ceremony.

The President arrives in top hat and tails, drawn in the imperial
coach and four. The blasts of the fanfare dance in the sunlight like bay-
onets. The dignitaries’ platform blooms with very beautiful noble-
women. The President reviews the platoon of cuirassiers. The ven-
dors circulate through the crowd with their badges and cockades.

Sun, smiles, prosperity, solemnity. On the other side of the fence, lost
in the throng of poor wretches, stand we, the tired old workers of the
tower. We look at each other, streams of tears running down our gray
beards. Ah, youth!


\end{document}
